\documentclass[11pt,a4paper]{article}
\usepackage{xeCJK}
\usepackage{fontspec}
\usepackage{xcolor}
\usepackage{fancyhdr}
\usepackage{booktabs}
\usepackage{array}
\usepackage{amssymb}
\usepackage{geometry}
\usepackage[protrusion=true]{microtype}
\geometry{margin=1.8cm,top=2cm,headheight=15pt}
\linespread{1.05}

\setmainfont{Baskerville}
\setCJKmainfont{Songti SC}
\newfontfamily{\starfont}{Songti SC}

\definecolor{wrongred}{RGB}{198,40,40}
\definecolor{wrongbg}{RGB}{253,235,233}
\definecolor{rulegray}{RGB}{200,200,200}

\setlength{\emergencystretch}{3em}
\setlength{\fboxsep}{3pt}
\setlength{\parindent}{0pt}

\newcommand{\checkmarkbox}{\fbox{\rule{0pt}{0.8ex}\rule{0.8ex}{0pt}}}
\newcommand{\STAR}{\textcolor{wrongred}{\starfont ★}}
% 题干里的填空线
\newcommand{\bk}{\rule[-0.55ex]{1.0em}{0.4pt}}
% 错答（红色）
\newcommand{\W}[1]{\textcolor{wrongred}{#1}}
% 正确答案（加粗）
\newcommand{\A}[1]{\textbf{#1}}
% 判定标记
\newcommand{\OK}{\textcolor{wrongred}{\ensuremath{\checkmark}}}
\newcommand{\NO}{\textcolor{wrongred}{$\times$}}
\newcommand{\Q}{\textcolor{wrongred}{\textbf{?}}}
\newcolumntype{L}[1]{>{\raggedright\arraybackslash}p{#1}}

\pagestyle{fancy}
\fancyhf{}
\fancyhead[L]{\footnotesize 英语语法 动词时态语态复习（错题订正）}
\fancyhead[R]{\footnotesize 赵文瀚}
\fancyfoot[C]{\thepage}
\renewcommand{\headrulewidth}{0.4pt}
\renewcommand{\footrulewidth}{0pt}

\begin{document}

\begin{center}
{\LARGE\bfseries 2026 学年高一英语国庆动词时态语态复习}\\[0.25em]
{\large 错题订正与复盘}\\[0.35em]
{\footnotesize 2029 届高一英语 · 时态语态专项（共 35 题）· 2026-10}
\end{center}

\vspace{-0.2em}
\noindent{\color{rulegray}\rule{\textwidth}{0.8pt}}

\vspace{0.5em}
\noindent\fcolorbox{black!15}{black!4}{%
  \parbox{\dimexpr\linewidth-2\fboxsep-2\fboxrule\relax}{%
    \small 日期：2026-10\qquad 来源：国庆动词时态语态复习（共 35 题）\qquad 姓名：赵文瀚\\[0.3em]
    判错：\textcolor{wrongred}{\textbf{10}}\,/\enspace 35\qquad
    判对：\textbf{24}\,/\enspace 35\qquad
    待核：\textbf{1}\,/\enspace 35\qquad
    重做日期：\underline{\hspace{2.4cm}}}}

\vspace{0.6em}
\noindent{\footnotesize 说明：下表按原卷题号列全 35 题。\W{红色}＝学生作答，\A{加粗}＝正确答案；\OK\ 判对、\NO\ 判错、\Q\ 待核。原卷\textbf{未被批改}（无红笔），本表对错为逐题重解所得。}

\vspace{0.4em}
\renewcommand{\arraystretch}{1.0}
\noindent{\footnotesize\begin{tabular}{@{}L{0.78cm}@{\hspace{0.08cm}}L{8.98cm}@{\hspace{0.08cm}}L{1.45cm}@{\hspace{0.08cm}}L{1.35cm}@{\hspace{0.08cm}}L{0.78cm}@{\hspace{0.08cm}}L{3.3cm}@{}}
\toprule
\textbf{题号} & \textbf{考点 · 题干要点} & \textbf{我的作答} & \textbf{正确答案} & \textbf{判定} & \textbf{错因} \\
\midrule
1  & 过去一连串动作：stepped into the office, \bk\ down and began & C & C & \OK & — \\
2  & She said she would telephone but we \bk\ from her so far & \W{B} & \A{A} & \NO & 现在完成时（so far） \\
3  & When I got to the cinema, the film \bk\ for ten minutes & C & C & \OK & — \\
4  & If city noises \bk\ from increasing, people \bk\ shout to be heard & A & A & \OK & — \\
5  & The last time I \bk\ Jane she \bk\ cotton in the fields & D & D & \OK & — \\
6  & he made it clear that he \bk\ office soon & \W{D} & \A{B} & \NO & 过去将来时（soon） \\
7  & The pen I \bk\ I \bk\ is on my desk, right under my nose & \W{A} & \A{B} & \NO & 时态呼应（thought…had lost） \\
8  & How long \bk\ each other before they \bk\ married? & D & D & \OK & — \\
9  & I don't really work here. I \bk\ until the new secretary arrives & \W{B} & \A{C} & \NO & 现在进行（临时状态） \\
10 & I need one more stamp before my collection \bk & D & D & \OK & — \\
11 & As she \bk\ the newspaper, Granny \bk\ asleep & B & B & \OK & — \\
12 & You don't need to describe her, I \bk\ her several times & B & B & \OK & — \\
13 & I don't think Jim saw me; he \bk\ into space & B & B & \OK & — \\
14 & — \bk\ my glasses? — Yes, I saw them on your bed a minute ago & D & D & \OK & — \\
15 & I first met Lisa three years ago. She \bk\ at a radio shop at the time & B & B & \OK & — \\
16 & — Who is Jerry Cooper? — \bk\ ? I saw you shaking hands with him（原卷印作 36.） & D & D & \OK & — \\
17 & Great changes \bk\ in my home town since liberation & \W{（涂改不清）} & \A{C} & \Q & 待核对原卷 \\
18 & The water will be further polluted unless some measures \bk & B & B & \OK & — \\
19 & He'll be an astronaut by the time he \bk\ thirty & A & A & \OK & — \\
20 & The \bk\ look on his face suggested that he \bk\ that & \W{A} & \A{B} & \NO & 分词 -ed/-ing；suggest＝表明 \\
21 & My brother is an actor. He \bk\ in several film so far & C & C & \OK & — \\
22 & I \bk\ ping-pong quite well, but I haven't had time to play since the new year & D & D & \OK & — \\
23 & Selecting a mobile phone is no easy task because technology \bk\ so rapidly & A & A & \OK & — \\
24 & I had just finished my work and \bk\ to take a shower & \W{B} & \A{D} & \NO & 过去进行（当时正要做） \\
25 & the city \bk\ in all directions in the past five years & \W{D} & \A{B} & \NO & 现在完成时（in the past five years） \\
26 & Never in all my life \bk\ so happy! & D & D & \OK & — \\
27 & You \bk\ television. Why not do something more active? & B & B & \OK & — \\
28 & My money \bk\ I must go to the bank before I've none in hand & \W{C} & \A{B} & \NO & 现在进行表将来；run out 不及物 \\
29 & — I shouldn't have been so rude. — You \bk\ your temper but that's OK & C & C & \OK & — \\
30 & recent science \bk\ that people who don't sleep well soon get ill & B & B & \OK & — \\
31 & a small box which \bk\ placed under the minister's car & C & C & \OK & — \\
32 & Because the shop \bk, all the T-shirts are sold at half price & \W{D} & \A{C} & \NO & 现在进行表"即将"（清仓） \\
33 & The crazy fans \bk\ patiently for two hours and they would wait till… & \W{C} & \A{B} & \NO & 过去完成进行时 \\
34 & Let's keep to the point or we \bk\ any decisions & A & A & \OK & — \\
35 & --Is this raincoat yours?  --No, mine \bk\ there behind the door & A & A & \OK & — \\
\bottomrule
\end{tabular}}

\vspace{0.3em}
\noindent{\footnotesize\textbf{原卷旁注}：（一）原卷\textbf{没有老师批改痕迹}（未见红笔），上表对错为逐题重解所得。（二）第 16 题原卷误印成「36.」；第 21 题 \emph{several film} 漏 \emph{s}；第 24 题选项 C 印作 \emph{C have started}（漏点）；第 16 题选项 C 印作 \emph{Didn't' you meet him yet}（多撇号）。（三）第 17 题学生答案为一个涂改墨团，认不出字母，正确答为 \emph{C}，请对照 {\xeCJKsetup{CJKecglue={}}\texttt{grammar\_01\_tense\_voice\_01\_原卷\_1.jpg}} 确认。（四）第 28 题 \emph{run out} 作「用完」是不及物动词，无被动式，故 C／D 均错；本表答案已与答案册逐题核对（第 7、24、32 题定为 B／D／C）。}

\newpage
\begin{center}
{\Large\bfseries 错因归类与复盘}
\end{center}
\vspace{-0.4em}
\noindent{\color{rulegray}\rule{\textwidth}{0.8pt}}

\vspace{0.5em}
\noindent 说明：按语法考点归类 10 条错题，供汇总对账；逐题二刷结果见第 1 页「判定」列。

\vspace{0.5em}
\renewcommand{\arraystretch}{1.35}
\begin{center}
\begin{tabular}{@{}p{6.6cm} p{7.6cm} c@{}}
\toprule
\textbf{易错考点} & \textbf{题号} & \textbf{题数} \\
\midrule
现在进行时（表"临时状态"或"即将发生"） & {\small 9, 28, 32} & 3 \\
现在完成时（与过去时／过去完成时分不清） & {\small 2, 25} & 2 \\
过去进行时（过去某时的背景动作） & {\small 24} & 1 \\
过去将来时（主句过去 + soon） & {\small 6} & 1 \\
过去完成进行时（截至过去某时"一直"在做） & {\small 33} & 1 \\
时态呼应（宾语从句：thought…had lost） & {\small 7} & 1 \\
分词与从句语气（-ed/-ing；suggest＝表明） & {\small 20} & 1 \\
\midrule
\textbf{合计} & \textbf{10 题} & \textbf{10} \\
\bottomrule
\end{tabular}
\end{center}

\vspace{0.5em}
\noindent\textbf{关联知识点标签}
\vspace{0.25em}

\noindent\begin{tabular}{@{}p{8.2cm}@{\hspace{0.4cm}}p{8.2cm}@{}}
\begin{minipage}[t]{\linewidth}
\small
\STAR\ \textbf{现在完成时信号词}：so far、in the past five years、since＋过去时间点\\[0.2em]
\STAR\ \textbf{过去完成进行时}：had been doing（截至过去某时一直在做）\\[0.2em]
\STAR\ \textbf{过去将来时}：主句是过去时 + soon／later $\to$ would do
\end{minipage} &
\begin{minipage}[t]{\linewidth}
\small
\STAR\ \textbf{现在进行时两用}：表"临时"（am just helping out）／表"即将"（is running out、is closing down）\\[0.2em]
\STAR\ \textbf{过去进行时}：when 从句提供背景，主句用过去进行（was starting to take a shower）\\[0.2em]
\STAR\ \textbf{宾语从句时态呼应}：I thought I had lost（主句过去 + 从句过去完成）\\[0.2em]
\STAR\ \textbf{分词作定语}：surprised look（人）／surprising news（物）\\[0.2em]
\STAR\ \textbf{suggest 双义}：建议（+ should do）／表明（用陈述语气）\\[0.2em]
\STAR\ \textbf{不及物动词无被动}：take place、run out、happen
\end{minipage} \\
\end{tabular}

\vspace{0.65em}
\noindent\fcolorbox{black!15}{black!4}{%
  \parbox{\dimexpr\linewidth-2\fboxsep-2\fboxrule\relax}{%
    \textbf{复盘记录}\\[0.35em]
    最薄弱的板块：\underline{\hspace{9.0cm}}\\[0.55em]
    主要错误原因：\checkmarkbox 时态混淆（完成时/进行时/过去将来）\quad \checkmarkbox 分词与词形\quad \checkmarkbox 主谓与语态\\[0.55em]
    本次复习的改进措施：\underline{\hspace{10.1cm}}\\[0.55em]
    \hspace{3.2em}\underline{\hspace{10.1cm}}\\[0.55em]
    下次复习日期：\underline{\hspace{2.2cm}}\qquad
    当前掌握度：\checkmarkbox 仍薄弱\quad \checkmarkbox 基本掌握\quad \checkmarkbox 已掌握
  }}

\end{document}
